% 编译方式：请使用 XeLaTeX 编译（xeCJK/ctex 中文支持依赖 XeTeX）
% xelatex 整体模型描述_LaTeX.tex
\documentclass[12pt,a4paper]{article}

% ==================== 中文支持 ====================
\usepackage[UTF8]{ctex}
% 注意：ctex 在 XeLaTeX 下会自动调用 xeCJK，无需手动 \usepackage{xeCJK}

% ==================== 页面布局 ====================
\usepackage[top=2.5cm, bottom=2.5cm, left=2.5cm, right=2.5cm]{geometry}
\usepackage{setspace}
\onehalfspacing

% ==================== 数学公式 ====================
\usepackage{amsmath,amssymb,amsfonts}
\usepackage{mathtools}

% ==================== 图表与颜色 ====================
\usepackage{graphicx}
\usepackage{float}
\usepackage{xcolor}
\usepackage{tikz}
\usetikzlibrary{positioning,arrows,shapes,fit,backgrounds,calc}

% ==================== 表格 ====================
\usepackage{booktabs}
\usepackage{multirow}
\usepackage{array}
\usepackage{tabularx}

% ==================== 算法伪代码 ====================
\usepackage{algorithm}
\usepackage{algorithmicx}
\usepackage{algpseudocode}

% ==================== 超链接与引用 ====================
\usepackage[hidelinks]{hyperref}
\usepackage[numbers,sort&compress]{natbib}

% ==================== 其他 ====================
\usepackage{enumitem}
\usepackage{caption}
\usepackage{subcaption}

% ==================== 自定义命令 ====================
\newcommand{\R}{\mathbb{R}}
\newcommand{\E}{\mathbb{E}}
\newcommand{\F}{\mathcal{F}}
\newcommand{\Sset}{\mathcal{S}}
\newcommand{\Aset}{\mathcal{A}}
\newcommand{\Cset}{\mathcal{C}}
\newcommand{\loss}{\mathcal{L}}
\newcommand{\reward}{\mathcal{R}}

\newcommand{\symptoms}{\mathbf{s}}
\newcommand{\selected}{\mathbf{c}}
\newcommand{\action}{\mathbf{a}}
\newcommand{\trueSE}{\mathbf{y}}
\newcommand{\predSE}{\hat{\mathbf{y}}}

\newcommand{\comment}[1]{\textcolor{gray}{#1}}

\title{\textbf{基于深度强化学习的中医辨证论治模型}}
\author{}
\date{\today}

\begin{document}
\maketitle

\section{问题定义}

中医辨证论治的核心任务可以形式化为一个\textbf{多标签分类}问题：给定一组刻下症（症状集合）$\symptoms = \{s_1, s_2, \ldots, s_n\}$，模型需要从中医学的证候要素（Syndrome Elements, SE）候选集合 $\Aset = \{a_1, a_2, \ldots, a_M\}$ 中，选择出与当前患者真实证候一致的子集 $\predSE \subseteq \Aset$。

与传统多标签分类不同，辨证论治具有以下特点：
\begin{enumerate}[label=(\arabic*)]
    \item \textbf{序列决策性}：医师在临床实践中并非一次性输出所有证候要素，而是逐步推理、逐条确认；
    \item \textbf{长尾分布}：证候要素在数据集中呈现严重的长尾分布，少数高频标签（如``气虚''）占据大多数样本，大量低频标签仅出现极少次数；
    \item \textbf{代价敏感性}：误诊（假阳性）和漏诊（假阴性）的代价不对称——在中医临床中，漏诊可能导致病情延误，其代价通常高于误诊；
    \item \textbf{集合级约束}：预测的证候要素数量应与真实数量接近，过多或过少均不合理。
\end{enumerate}

为此，我们将辨证论治建模为一个\textbf{马尔可夫决策过程（MDP）}，并采用\textbf{深度Q网络（DQN）}进行求解。

\section{马尔可夫决策过程建模}

\subsection{状态空间 $\Sset$}

状态向量由两部分拼接而成，均为 multi-hot 编码：

\begin{equation}
    \mathbf{s}_t = \big[\,\underbrace{\mathbf{s}^{\text{symp}}}_{n\text{维}}\, \|\, \underbrace{\mathbf{s}^{\text{selected}}_t}_{M\text{维}}\,\big] \in \{0,1\}^{n+M}
\end{equation}

其中：
\begin{itemize}
    \item $\mathbf{s}^{\text{symp}}$：当前患者的症状 multi-hot 向量，若症状 $i$ 出现则为 1；
    \item $\mathbf{s}^{\text{selected}}_t$：截至第 $t$ 步已选择的证候要素 multi-hot 向量。
\end{itemize}

初始状态 $\mathbf{s}_0$ 的已选部分全为零向量。

\subsection{动作空间 $\Aset$}

动作空间包含两类动作：
\begin{itemize}
    \item \textbf{标签动作}：选择第 $j$ 个证候要素，$j \in \{0, 1, \ldots, M-1\}$；
    \item \textbf{停止动作}：$a_{\text{stop}} = M$，表示终止当前病例的辨证过程。
\end{itemize}

合法动作需满足：
\begin{itemize}
    \item 已选标签不可重复选择（通过 Q 值 mask 实现）；
    \item 停止动作在至少选择了 $\text{min\_actions}$ 个标签后才可选。
\end{itemize}

\subsection{状态转移}

选择标签动作 $a_t = j$（$j \neq \text{stop}$）后：
\begin{equation}
    \mathbf{s}_{t+1}^{\text{selected}} = \mathbf{s}_t^{\text{selected}} \text{ 且将第 } j \text{ 位置为 } 1
\end{equation}

选择停止动作后，episode 终止。

\subsection{奖励函数}

奖励函数是本模型的核心设计之一。我们支持四种模式，默认使用 \textbf{tail\_cost\_curiosity} 模式。该模式的奖励函数设计包含以下组件：

\subsubsection{势函数塑形}

为了缓解稀疏奖励问题，我们引入基于势函数的奖励塑形（Potential-Based Reward Shaping）。定义势函数 $\Phi(\selected)$ 来衡量当前已选集合 $\selected$ 的质量：

\begin{equation}
\begin{aligned}
\Phi(\selected) = &\;\alpha_{\Phi} \cdot \Big(f_{\text{log}}(F_{\beta}(\selected)) - \Phi_{\text{baseline}}\Big) \\
& - \lambda_{\text{card}} \cdot \frac{|\selected| - |\trueSE|}{|\trueSE|} \\
& - \lambda_{\text{fn}} \cdot \frac{w_{\text{FN}}(\selected)}{w_{\text{true}}} \\
& + \lambda_{\text{curiosity}} \cdot \frac{1}{w_{\text{true}}} \sum_{a \in \selected \cap \trueSE} h(a) \cdot w(a)
\end{aligned}
\end{equation}

其中：
\begin{itemize}
    \item $F_{\beta}(\selected)$ 是加权平衡 F-score，$\beta = 1.30$，增强召回率权重；
    \item $f_{\text{log}}(x) = \frac{\ln(1 + g \cdot x)}{\ln(1 + g)}$ 是对数边际效用递减变换（$g = 9.0$），使得低 F1 区域的改进获得更大奖励增量；
    \item $w(a)$ 是标签稀有度权重（详见第\ref{sec:weight}节）；
    \item $h(a)$ 是预训练好奇心提示分数（详见第\ref{sec:curiosity}节）；
    \item $w_{\text{FN}}$ 和 $w_{\text{true}}$ 分别为漏选加权和与真实标签加权和。
\end{itemize}

步骤奖励为：
\begin{equation}
    \reward_t = \gamma \cdot \Phi(\selected_{t+1}) - \Phi(\selected_t) - \lambda_{\text{step}} + \reward_t^{\text{bonus}} - \reward_t^{\text{penalty}}
\end{equation}

\subsubsection{步骤奖励的额外项}

\begin{equation}
\reward_t^{\text{bonus}} = 
\begin{cases}
    \lambda_{\text{rare}} \cdot \max(0, w(a_t)-1) & \text{若 } a_t \in \trueSE \text{（稀有标签奖励）} \\
    + \; \reward_t^{\text{curiosity}} & \text{（好奇心奖励）} \\
    + \; \lambda_{\text{explore}} \cdot h(a_t) \cdot \max(0, w(a_t)-1) \cdot (1-p) & \text{（探索奖励，随训练衰减）} \\[6pt]
    0 & \text{否则}
\end{cases}
\end{equation}

\begin{equation}
\reward_t^{\text{penalty}} = 
\begin{cases}
    \lambda_{\text{fp}}(p) \cdot w_{\text{fp}}(a_t) & \text{若 } a_t \notin \trueSE \text{（假阳性惩罚）} \\
    + \; \lambda_{\text{over}} \cdot \big(\max(0, |\selected_{t+1}| - |\trueSE|)\big)^{\rho} & \text{（二次多选惩罚，}\rho = 2.0\text{）}
\end{cases}
\end{equation}

其中假阳性惩罚系数 $\lambda_{\text{fp}}(p)$ 随训练进度 $p$ 线性插值：从 $\lambda_{\text{fp}}^{\text{start}} = 0.12$ 到 $\lambda_{\text{fp}}^{\text{end}} = 0.36$。

\subsubsection{终端奖励}

当智能体选择停止动作时，给予终端奖励：
\begin{equation}
\begin{aligned}
\reward_{\text{terminal}} = &\;\lambda_{\text{F1}} \cdot \big(\omega_s \cdot \text{F1}_{\text{sample}} + \omega_b \cdot F_{\beta}(\selected)\big) \\
& - \lambda_{\text{fn}}^{\text{stop}} \cdot \frac{w_{\text{FN}}}{w_{\text{true}}} - \lambda_{\text{fp}}^{\text{stop}} \cdot w_{\text{FP}} \\
& - \lambda_{\text{card}}^{\text{term}} \cdot \big||\selected| - |\trueSE|\big| \\
& + \mathbb{I}[\text{完全匹配}] \cdot \lambda_{\text{exact}} \\
& + \lambda_{\text{tail}} \cdot \frac{\sum_{a \in \selected \cap \trueSE} \max(0, w(a)-1)}{\sum_{a \in \trueSE} \max(0, w(a)-1)} \\
& - \lambda_{\text{miss}} \cdot \frac{1}{w_{\text{true}}} \sum_{a \in \trueSE \setminus \selected} h(a) \cdot w(a) \\
& - \lambda_{\text{under}} \cdot \max(0, |\trueSE| - |\selected|)
\end{aligned}
\end{equation}

\subsubsection{关键超参数汇总}

\begin{table}[H]
\centering
\caption{tail\_cost\_curiosity 模式奖励函数超参数}
\begin{tabular}{lcl}
\toprule
\textbf{参数} & \textbf{符号} & \textbf{默认值} \\
\midrule
势函数缩放系数 & $\alpha_{\Phi}$ & 1.60 \\
对数增益系数 & $g$ & 9.0 \\
势函数基线 & $\Phi_{\text{baseline}}$ & 0.0 \\
步惩罚 & $\lambda_{\text{step}}$ & 0.008 \\
多选惩罚（起始/终止） & $\lambda_{\text{over}}$ & 0.08 \\
多选惩罚幂次 & $\rho$ & 2.0 \\
假阳性惩罚（起始/终止） & $\lambda_{\text{fp}}^{\text{start}}/\lambda_{\text{fp}}^{\text{end}}$ & 0.12/0.36 \\
稀有标签奖励 & $\lambda_{\text{rare}}$ & 0.40 \\
数量误差惩罚 & $\lambda_{\text{card}}$ & 0.03 \\
漏选惩罚 & $\lambda_{\text{fn}}$ & 0.34 \\
终端 F1 缩放 & $\lambda_{\text{F1}}$ & 2.80 \\
样本 F1 权重 & $\omega_s$ & 0.20 \\
平衡 F1 权重 & $\omega_b$ & 0.80 \\
完全匹配奖励 & $\lambda_{\text{exact}}$ & 0.60 \\
终端漏选惩罚 & $\lambda_{\text{fn}}^{\text{stop}}$ & 0.28 \\
终端假阳性惩罚 & $\lambda_{\text{fp}}^{\text{stop}}$ & 0.30 \\
尾部召回奖励 & $\lambda_{\text{tail}}$ & 0.45 \\
预训练漏选惩罚 & $\lambda_{\text{miss}}$ & 0.24 \\
不确定性探索奖励 & $\lambda_{\text{uncertainty}}$ & 0.25 \\
预训练漏检奖励 & $\lambda_{\text{pretrain\_miss}}$ & 0.35 \\
好奇心势能缩放 & $\lambda_{\text{curiosity}}$ & 0.18 \\
探索奖励 & $\lambda_{\text{explore}}$ & 0.15 \\
终端漏选数量惩罚 & $\lambda_{\text{under}}$ & 0.12 \\
假阳性权重缩放 & $\lambda_{\text{fp\_scale}}$ & 0.35 \\
\bottomrule
\end{tabular}
\end{table}

\subsection{标签稀有度权重设计}
\label{sec:weight}

为缓解长尾分布问题，我们为每个证候要素 $a$ 分配一个稀有度权重 $w(a)$。该权重通过对数压缩方法计算：

\begin{equation}
    w(a) = \min\Big(w_{\max},\; \max\big(w_{\min},\; w_{\min} + \lambda_{\log} \cdot \ln\big(\frac{\max_c \text{count}(c)}{\text{count}(a)}\big)\big)\Big)
\end{equation}

其中 $\text{count}(a)$ 为标签 $a$ 在训练集中的出现次数。对数压缩相比幂函数压缩（$(\max_c / \text{count}(a))^\gamma$）能有效避免极端长尾标签（如仅出现 2 次的标签）的权重达到数百倍的不合理范围。当前参数：$w_{\min}=0.6$，$w_{\max}=6.0$，$\lambda_{\log}=0.65$。

此外，为了缓解低支持度标签的权重估计不确定性，我们引入支持度置信度因子：
\begin{equation}
    \text{conf}(a) = c_{\min} + (1 - c_{\min}) \cdot \frac{\text{support}(a)}{\text{support}(a) + k}
\end{equation}
其中 $c_{\min}=0.35$，$k=5.0$。实际使用的有效权重为：
\begin{equation}
    w_{\text{eff}}(a) = 1 + (w(a) - 1) \cdot \text{conf}(a)
\end{equation}

\subsection{好奇心机制}
\label{sec:curiosity}

为解决预训练模型对低频标签预测置信度低的问题，我们引入了基于预训练不确定性的好奇心机制。

\subsubsection{预训练探索提示}

在监督预训练之后，使用预训练模型对训练集中的每个症状集合进行推理，为每个候选标签计算提示分数：
\begin{equation}
    h(a) = \max\big(0,\; \min\big(1,\; 0.55 \cdot (1 - p(a)) + 0.45 \cdot u(a)\big)\big)
\end{equation}
其中 $p(a)$ 是预训练模型对标签 $a$ 的预测概率，$u(a) = 1 - |2p(a) - 1|$ 是预测不确定性（当 $p(a)=0.5$ 时最大）。

同时记录``漏检标记'' $m(a)$：当标签 $a$ 在预训练模型的预测概率排序中超出 $\max(|\trueSE|+2, 5)$ 位时标记为漏检。

好奇心奖励在智能体选中真实标签时激活：
\begin{equation}
    \reward_t^{\text{curiosity}} = \lambda_{\text{uncertainty}} \cdot h(a_t) \cdot w_{\text{eff}}(a_t) \;+\; \lambda_{\text{pretrain\_miss}} \cdot m(a_t) \cdot w_{\text{eff}}(a_t)
\end{equation}

\subsubsection{周期性更新}

由于 RL 训练过程中网络权重持续变化，预训练提示会逐渐过时。我们每 $\text{refresh\_interval}=10$ 次迭代重新计算一次提示分数，确保好奇心机制始终反映当前网络的不确定性。

\subsection{课程学习}

训练采用三阶段课程学习策略：
\begin{itemize}
    \item \textbf{第一阶段}（前 30\% 迭代）：仅使用真实标签数 $\leq 2$ 的简单样本；
    \item \textbf{第二阶段}（30\%--60\% 迭代）：扩展到真实标签数 $\leq 4$ 的样本；
    \item \textbf{第三阶段}（后 40\% 迭代）：使用全部训练样本。
\end{itemize}

\section{深度Q网络架构}

\subsection{网络结构：Set-Aware Dueling DQN with NoisyNet}

我们采用 Set-Aware Dueling DQN 作为默认网络架构。该架构的核心创新在于将状态向量的两部分——症状向量和已选证候要素向量——分别编码，再融合交互信息：

\begin{figure}[H]
\centering
\begin{tikzpicture}[
    node distance=0.8cm,
    block/.style={rectangle, draw, minimum width=3cm, minimum height=0.7cm, align=center, rounded corners=2pt},
    smallblock/.style={rectangle, draw, minimum width=2cm, minimum height=0.5cm, align=center, rounded corners=2pt},
    arrow/.style={->, >=stealth, thick}
]
    % 输入
    \node[block, fill=blue!10] (symptom_input) {症状向量 $\symptoms$ \\ ($n$ 维 multi-hot)};
    \node[block, fill=blue!10, right=1cm of symptom_input] (selected_input) {已选标签向量 $\selected_t$ \\ ($M$ 维 multi-hot)};

    % 编码器
    \node[block, fill=green!10, below=of symptom_input] (symptom_encoder) {症状编码器 \\ Linear($n$, 256) + LayerNorm + ReLU \\ + Dropout(0.15) + ResidualBlock};
    \node[block, fill=green!10, below=of selected_input] (selected_encoder) {已选编码器 \\ Linear($M$, 256) + LayerNorm + ReLU \\ + Dropout(0.15) + ResidualBlock};

    % 融合层
    \node[block, fill=orange!10, below=1.5cm of $(symptom_encoder.south)!0.5!(selected_encoder.south)$] (joint) {联合融合层 \\ Concat(症状特征, 已选特征, 逐元素积) \\ Linear(768, 256) + ResidualBlock \\ + Linear(256, 128)};

    % Dueling 分支
    \node[block, fill=red!10, below left=0.8cm and 1cm of joint] (value) {价值流 Value Stream \\ Linear(128, 128) \\ + Linear(128, 1)};
    \node[block, fill=red!10, below right=0.8cm and 1cm of joint] (advantage) {优势流 Advantage Stream \\ Linear(128, 128) \\ + Linear(128, $M+1$)};

    % 输出
    \node[block, fill=gray!10, below=1.2cm of $(value.south)!0.5!(advantage.south)$] (output) {Q 值输出 \\ $Q(s, a) = V(s) + A(s, a) - \frac{1}{|\Aset|}\sum_{a'}A(s, a')$};

    % 箭头
    \draw[arrow] (symptom_input) -- (symptom_encoder);
    \draw[arrow] (selected_input) -- (selected_encoder);
    \draw[arrow] (symptom_encoder) -- (joint);
    \draw[arrow] (selected_encoder) -- (joint);
    \draw[arrow] (joint) -- (value);
    \draw[arrow] (joint) -- (advantage);
    \draw[arrow] (value) -- (output);
    \draw[arrow] (advantage) -- (output);

\end{tikzpicture}
\caption{Set-Aware Dueling DQN 网络架构}
\label{fig:network}
\end{figure}

\subsection{NoisyNet 探索}

在融合层中，我们使用 Factorized NoisyNet 线性层替代标准线性层，实现状态依赖的自适应探索。NoisyNet 将权重和偏置参数化为：
\begin{equation}
    \mathbf{W} = \boldsymbol{\mu}_w + \boldsymbol{\sigma}_w \odot \boldsymbol{\varepsilon}_w, \quad \mathbf{b} = \boldsymbol{\mu}_b + \boldsymbol{\sigma}_b \odot \boldsymbol{\varepsilon}_b
\end{equation}
其中 $\boldsymbol{\mu}$ 是可学习参数，$\boldsymbol{\sigma}$ 是可学习的噪声尺度，$\boldsymbol{\varepsilon}$ 是每次前向传播时采样的零均值噪声（Factorized Gaussian）。

\subsection{模型变体}

除默认的 Set-Aware Dueling DQN 外，我们还实现了以下变体用于消融实验：
\begin{itemize}
    \item \textbf{DQN (Legacy)}：两层 MLP，无 LayerNorm；
    \item \textbf{DQN}：两层 MLP + LayerNorm + Dropout；
    \item \textbf{Dueling DQN}：标准 Dueling 架构，共享特征层 + 分离的价值流和优势流；
    \item \textbf{Residual Dueling DQN}：在特征提取层加入两个残差块；
    \item \textbf{Set-Aware Dueling DQN}：分路编码症状和已选标签，加入交互项；
    \item \textbf{Set-Aware Dueling DQN + NoisyNet}：默认模型，融合层使用 NoisyNet。
\end{itemize}

\section{训练流程}

\subsection{监督预训练}

在 RL 训练之前，我们首先进行少量 epoch（默认 5 个）的监督预训练，使用 BCEWithLogitsLoss 损失函数。预训练数据通过枚举每条训练样本的所有标签排列顺序生成中间状态：

\begin{algorithm}[H]
\caption{监督预训练数据生成}
\begin{algorithmic}[1]
\For{each 训练样本 $(\symptoms, \trueSE)$}
    \State $\text{true\_actions} \gets \text{映射到动作索引}$
    \State $\text{trajectories} \gets [\text{true\_actions}]$
    \For{$k = 2$ to $\text{pretrain\_all\_permutations}$}
        \State $\text{shuffled} \gets \text{随机打乱}(\text{true\_actions})$
        \State $\text{trajectories}.\text{append}(\text{shuffled})$
    \EndFor
    \For{each trajectory in trajectories}
        \For{$t = 0$ to $|\text{trajectory}|$}
            \State $\selected_t \gets \text{trajectory}[:t]$
            \State $\text{state} \gets \text{build\_state\_vector}(\symptoms, \selected_t)$
            \State $\text{target}[a] \gets 1$ for $a \in \trueSE \setminus \selected_t$
            \State $\text{target}[\text{stop}] \gets 1$ if $\selected_t = \trueSE$ else 0
            \State 保存 $(\text{state}, \text{target})$
        \EndFor
    \EndFor
\EndFor
\end{algorithmic}
\end{algorithm}

预训练使用 AdamW 优化器，学习率 0.001，batch size 256，正样本权重按逆频次计算并截断到 $[1, 18]$。

\subsection{经验池预热}

预训练完成后，使用专家轨迹（真实标签排列）与环境交互，将产生的 transitions 存入经验回放池。每条训练样本生成 $\text{pretrain\_all\_permutations}$ 条（默认 2 条）轨迹。

\subsection{主训练循环}

\begin{algorithm}[H]
\caption{DQN 主训练循环（单次迭代）}
\begin{algorithmic}[1]
\For{each episode $i = 0$ to $\text{num\_episodes}-1$}
    \State 更新训练进度 $p \gets i / (\text{num\_episodes}-1)$
    \State 根据课程学习策略筛选当前 episode 的训练样本
    \For{each 训练样本}
        \State 重置环境，获得初始状态 $\mathbf{s}_0$
        \State $\selected \gets []$
        \Repeat
            \State $\epsilon$-greedy 选择动作 $a_t$
            \State 环境执行 $a_t$，获得 $(\mathbf{s}_{t+1}, r_t, \text{done})$
            \State 记录 transition 到缓冲区
            \If{$a_t \neq \text{stop}$}
                \State $\selected.\text{append}(a_t)$
            \EndIf
        \Until{done}
        \State 将缓冲区中的 transitions 以 n-step 方式存入 ReplayMemory
        \While{$|\text{ReplayMemory}| \geq \text{batch\_size}$ \textbf{and} 累计步数 $\geq$ optimize\_interval}
            \State 从 ReplayMemory 采样 batch
            \State 计算 Double DQN 目标值
            \State 计算 TD 损失（SmoothL1Loss）+ 辅助监督损失 + 锚点损失
            \State 梯度裁剪 + 优化器更新
            \State 软更新目标网络：$\theta^- \gets (1-\tau)\theta^- + \tau\theta$
            \State 更新 PER 优先级
        \EndWhile
    \EndFor
    \If{$(i+1) \bmod \text{refresh\_interval} = 0$}
        \State 更新预训练探索提示 $\{h(a)\}$
    \EndIf
\EndFor
\end{algorithmic}
\end{algorithm}

\subsection{关键训练技术}

\begin{enumerate}[label=(\arabic*)]
    \item \textbf{Double DQN}：使用 policy 网络选择动作，target 网络评估价值，缓解过估计；
    \item \textbf{Prioritized Experience Replay (PER)}：使用 SumTree 实现 $O(\log N)$ 的优先级采样，采样概率 $P(i) \propto (|\delta_i| + \epsilon)^\alpha$（$\alpha=0.6$），重要性采样权重 $\beta$ 从 0.4 线性退火到 1.0；
    \item \textbf{n-step Return}：使用 3-step bootstrapping 加速价值传播；
    \item \textbf{辅助监督损失}：在 RL 训练中持续施加辅助 BCE 损失，权重从 0.15 线性衰减到 0.03，防止灾难性遗忘；
    \item \textbf{预训练锚点损失}：对预训练后的网络参数施加 L2 正则化，防止 RL 训练过度偏离预训练知识（当前权重为 0，未启用）；
    \item \textbf{软目标网络更新}：$\tau = 0.01$，每次优化步骤后更新；
    \item \textbf{梯度裁剪}：最大梯度范数 5.0；
    \item \textbf{PER 优先级增强}：存入 transition 时，稀有标签额外增加 $w(a) \cdot \lambda_{\text{rare\_priority}}$ 的优先级，预训练漏检标签额外增加 $h(a) \cdot \lambda_{\text{miss\_priority}}$ 的优先级。
\end{enumerate}

\subsection{关键超参数汇总}

\begin{table}[H]
\centering
\caption{DQN 训练超参数}
\begin{tabular}{lcl}
\toprule
\textbf{参数} & \textbf{符号} & \textbf{默认值} \\
\midrule
训练迭代数 & $\text{num\_episodes}$ & 50 \\
批次大小 & $\text{batch\_size}$ & 64 \\
经验池容量 & $\text{mem\_capacity}$ & 10000 \\
折扣因子 & $\gamma$ & 0.95 \\
探索率（起始/终止/衰减） & $\epsilon_s / \epsilon_e / \epsilon_d$ & 0.7 / 0.02 / 10000 \\
软更新系数 & $\tau$ & 0.01 \\
学习率 & $\eta$ & $3 \times 10^{-4}$ \\
权重衰减 & $\lambda_{\text{wd}}$ & $10^{-4}$ \\
优化器 & - & AdamW (amsgrad=True) \\
隐藏层1大小 & $n_{\text{units}}$ & 256 \\
隐藏层2大小 & $n_{\text{units2}}$ & 128 \\
Dropout 率 & - & 0.15 \\
预训练 epoch & - & 5 \\
预训练学习率 & - & 0.001 \\
预训练标签排列数 & - & 2 \\
辅助监督权重（起始/终止） & - & 0.15 / 0.03 \\
PER $\alpha$ & - & 0.6 \\
PER $\beta$ 起始 & - & 0.4 \\
PER $\beta$ 退火帧数 & - & 5000 \\
n-step & - & 3 \\
优化间隔 & - & 5 \\
课程学习阶段1/2 & - & 0.3 / 0.6 \\
最少选择标签数 & $\text{min\_actions}$ & 2 \\
稀有标签优先级缩放 & - & 1.1 \\
预训练漏检优先级缩放 & - & 0.7 \\
提示更新间隔 & - & 10 \\
\bottomrule
\end{tabular}
\end{table}

\section{评估指标}

\subsection{核心指标}

\begin{enumerate}[label=(\arabic*)]
    \item \textbf{样本 F1（Sample F1）}：对每个病例先计算 Precision 和 Recall，再算 F1，最后对所有病例取平均。这是本任务最直观的核心指标。
    \begin{equation}
        \text{Sample F1} = \frac{1}{N} \sum_{i=1}^{N} \frac{2 \cdot P_i \cdot R_i}{P_i + R_i}
    \end{equation}
    \item \textbf{严格准确率（Exact Match）}：预测集合与真实集合完全一致的病例比例。
    \item \textbf{Macro F1}：先计算每个标签的 F1，再对所有标签取平均，更能反映低频标签的表现。
    \item \textbf{Micro F1}：将所有标签的 TP/FP/FN 汇总后计算全局 F1，更受高频标签影响。
\end{enumerate}

\subsection{辅助指标}

\begin{enumerate}[label=(\arabic*)]
    \item \textbf{样本 Jaccard}：预测集合与真实集合的交并比；
    \item \textbf{标签级 Accuracy / Hamming Loss}：全局标签级正确率/错误率；
    \item \textbf{Macro Precision / Recall}：标签级精确率和召回率；
    \item \textbf{Supported Macro F1}：仅对训练集中出现过的标签计算 Macro F1；
    \item \textbf{Hit Rate}：至少命中一个真实标签的病例比例；
    \item \textbf{Empty Prediction Rate}：未推荐任何标签的病例比例；
    \item \textbf{勿选率（False Selection Rate）}：预测集合中假阳性标签的占比；
    \item \textbf{漏选率（Miss Selection Rate）}：真实集合中未被预测的标签占比；
    \item \textbf{平均推荐数 / 平均真实数}：观察模型推荐数量是否合理；
    \item \textbf{平均多选数 / 平均漏选数}：量化数量偏差；
    \item \textbf{逐标签 P/R/F1/Support/PredCount}：定位每个标签的预测表现。
\end{enumerate}

\section{消融实验框架}

消融实验通过修改命令行参数来系统地控制单一变量，以验证各模块的独立贡献。主要消融维度包括：

\begin{itemize}
    \item \textbf{奖励模式消融}：对比 legacy、potential\_f1、macro\_balanced、tail\_cost\_curiosity 四种奖励模式；
    \item \textbf{网络架构消融}：对比六种 Q 网络架构；
    \item \textbf{核心组件消融}：通过参数开关控制预训练、PER、辅助监督、课程学习、n-step、NoisyNet 等模块的启用与否；
    \item \textbf{奖励函数子模块消融}：控制好奇心奖励、探索奖励、对数 F1 变换、二次惩罚等子模块的开关。
\end{itemize}

所有消融实验统一控制随机种子，使用相同的训练/测试集划分，确保结果可比。

\section{模型文件结构}

\begin{table}[H]
\centering
\caption{项目文件结构}
\begin{tabular}{ll}
\toprule
\textbf{文件} & \textbf{功能} \\
\midrule
\texttt{main.py} & 主入口，解析参数，协调训练/评估流程 \\
\texttt{env.py} & RL 环境定义，含状态空间、动作空间、奖励函数 \\
\texttt{trainer.py} & DQN 训练器，含预训练、经验池预热、主训练循环 \\
\texttt{model.py} & 六种 Q 网络架构定义 \\
\texttt{data.py} & 数据加载、划分、标签权重计算 \\
\texttt{state.py} & 状态向量构建与预测动作序列生成 \\
\texttt{action.py} & 动作选择（epsilon-greedy + mask） \\
\texttt{memory.py} & 经验回放池（含 SumTree PER） \\
\texttt{evaluation.py} & 评估指标计算 \\
\texttt{metrics.py} & 底层集合指标（TP/FP/FN/F1/Jaccard） \\
\texttt{constants.py} & 全局常量 \\
\texttt{utils.py} & 工具函数（logger、seed） \\
\texttt{predict\_core.py} & 独立预测共享核心 \\
\texttt{predict\_trace.py} & 逐步辨证预测入口 \\
\texttt{plot\_training.py} & 训练日志可视化 \\
\bottomrule
\end{tabular}
\end{table}

\section{技术创新点总结}

\begin{enumerate}[label=(\arabic*)]
    \item \textbf{序列化辨证建模}：将中医辨证论治建模为 MDP，通过逐步选择证候要素模拟临床推理过程；
    \item \textbf{非线性奖励函数设计}：
    \begin{itemize}
        \item 对数边际效用递减的 F1 变换，使低 F1 区域梯度更大；
        \item 二次多选惩罚，少量容忍、大量重罚；
        \item 对数压缩的标签稀有度权重，避免极端长尾值主导训练；
        \item 支持度置信度因子，缓解低支持度标签的权重估计不确定性。
    \end{itemize}
    \item \textbf{好奇心驱动的尾部探索}：通过预训练不确定性提示 + 周期性更新机制，引导模型关注低频标签；
    \item \textbf{Set-Aware Dueling 架构}：分路编码症状证据和已选集合，并通过逐元素积显式建模二者交互；
    \item \textbf{NoisyNet 自适应探索}：替代传统 epsilon-greedy，实现状态依赖的探索策略；
    \item \textbf{多阶段课程学习}：从简单到复杂逐步引入训练样本，降低学习难度；
    \item \textbf{综合训练技术}：融合 Double DQN、PER、n-step Return、辅助监督损失、梯度裁剪等多项稳定训练技术。
\end{enumerate}

\end{document}
